% Part: computability
% Chapter: recursive-functions
% Section: introduction

\documentclass[../../../include/open-logic-section]{subfiles}

\begin{document}

\olfileid{cmp}{rec}{int}
\olsection{ভূমিকা}

গণনসাধ্যতার একটি গাণিতিক তত্ত্ব গড়ে তুলতে হলে প্রথমেই গণনসাধ্যতার
একটি \emph{মডেল} গড়ে তুলতে হয়। এখন আমরা গণনসাধ্যতাকে কম্পিউটার যে
ধরনের কাজ করে বলে ভাবি, আর কম্পিউটার সংকেত নিয়ে কাজ করে। কিন্তু
গণনসাধ্যতার তত্ত্বগুলির বিকাশের শুরুতে গণনার আদর্শ উদাহরণ ছিল
\emph{সংখ্যাগত} গণনা। গণিতবিদেরা সংখ্যাতাত্ত্বিক অপেক্ষক নিয়ে সবসময়ই
আগ্রহী ছিলেন, অর্থাৎ এমন অপেক্ষক $f\colon \Nat^n \to \Nat$ নিয়ে, যা
গণনা করা যায়। তাই গণনসাধ্যতার তত্ত্বের শুরুতে এই ধরনের অপেক্ষকগুলিই
অধ্যয়ন করা হয়েছিল, এতে আশ্চর্যের কিছু নেই। গণনসাধ্য সংখ্যাগত অপেক্ষকের
সবচেয়ে পরিচিত উদাহরণ—যেমন স্বাভাবিক সংখ্যার যোগ, গুণ ও ঘাত—একটি
আকর্ষণীয় বৈশিষ্ট্য ভাগ করে: এগুলিকে \emph{পুনরাবৃত্তভাবে} সংজ্ঞায়িত
করা যায়। সুতরাং পুনরাবৃত্ত সংজ্ঞার ভিত্তিতে \emph{গণনসাধ্য অপেক্ষক}-এর
একটি সাধারণ সংজ্ঞা দেওয়ার চেষ্টা করাই স্বাভাবিক। সংখ্যাতাত্ত্বিক অপেক্ষক
পুনরাবৃত্তভাবে সংজ্ঞায়িত করার বহু সম্ভাব্য পদ্ধতির মধ্যে একটি বিশেষ
সরল বিন্যাস এখানে কেন্দ্রীয় হয়ে উঠবে: তথাকথিত \emph{আদিম পুনরাবৃত্তি}।

গণনসাধ্য অপেক্ষকের পাশাপাশি গণনসাধ্য সেট ও সম্বন্ধ নিয়েও আমাদের আগ্রহ
থাকতে পারে। কোনো সেট গণনসাধ্য, যদি প্রদত্ত কোনো সংখ্যা সেই সেটের !!a{element}
কি না, তা আমরা গণনা করে নির্ণয় করতে পারি; এবং কোনো সম্বন্ধ গণনসাধ্য,
যদি কোনো টাপল $\tuple{n_1, \dots, n_k}$ সেই সম্বন্ধের !!a{element} কি না,
তা গণনা করে নির্ণয় করতে পারি। কোনো সেট বা সম্বন্ধের \emph{চরিত্রসূচক
অপেক্ষক} বিবেচনা করলে গণনসাধ্য সেট ও সম্বন্ধের আলোচনা গণনসাধ্য অপেক্ষকের
আলোচনার অন্তর্ভুক্ত করা যায়। তাই আদিম পুনরাবৃত্ত সম্বন্ধও সংজ্ঞায়িত
করতে পারি; যেমন, “$n$ দিয়ে $m$ নিঃশেষে বিভাজ্য” সম্বন্ধটি আদিম
পুনরাবৃত্ত সম্বন্ধ।

আদিম পুনরাবৃত্ত অপেক্ষক—অর্থাৎ কেবল আদিম পুনরাবৃত্তি ব্যবহার করে
সংজ্ঞায়িত অপেক্ষক—তবে একমাত্র গণনসাধ্য সংখ্যাতাত্ত্বিক অপেক্ষক নয়।
আদিম পুনরাবৃত্তির বহু সাধারণীকরণ বিবেচনা করা হয়েছে, কিন্তু অপেক্ষক
গণনার সবচেয়ে শক্তিশালী ও বহুলস্বীকৃত অতিরিক্ত পদ্ধতি হলো সীমাহীন
অনুসন্ধান। এর ফলে \emph{আংশিক পুনরাবৃত্ত অপেক্ষক}-এর সংজ্ঞা এবং তার
সঙ্গে সম্পর্কিত \emph{সাধারণ পুনরাবৃত্ত অপেক্ষক}-এর সংজ্ঞা পাওয়া যায়।
সাধারণ পুনরাবৃত্ত অপেক্ষক গণনসাধ্য ও সর্বত্র-সংজ্ঞায়িত; এই সংজ্ঞা ঠিক
সেই আংশিক পুনরাবৃত্ত অপেক্ষকগুলিকে চিহ্নিত করে, যেগুলি সর্বত্র-সংজ্ঞায়িত।
পুনরাবৃত্ত অপেক্ষক অন্য সব গণনামডেল (টুরিং যন্ত্র, ল্যামডা ক্যালকুলাস
ইত্যাদি) অনুকরণ করতে পারে, তাই এগুলি গণনার স্বীকৃত বহু মডেলের একটি
প্রতিনিধিত্ব করে।


\end{document}
